\RequirePackage{plautopatch}
\documentclass[platex、dvipdfmx、a4paper]{jsarticle}

\usepackage{amsmath}
\usepackage{amssymb}
\newcommand{\hb}[2]{\frac{\sectionial #1}{\sectionial #2}}
\newcommand{\bbn}[2]{\frac{d #1}{d #2}}
\newcommand{\hhb}[2]{\frac{\sectionial^2 #1}{\sectionial #2^2}}
\newcommand{\bbbn}[2]{\frac{d^2 #1}{d #2^2}}
\newcommand{\vectxy}[2]{\begin{pmatrix}#1\\#2\end{pmatrix}}
\newcommand{\vectxyz}[3]{\begin{pmatrix}#1\\#2\\#3\end{pmatrix}}
\newcommand{\vect}[1]{\boldsymbol{\mathrm{#1}}}

\newcommand{\ssum}[3]{\sum_{#1=#2}^#3}

\title{物理数学II演習レポート2}
\author{嶋田 悠音}
\date{2026年6月22日}

\begin{document}
  \maketitle
  \subsection*{問題3$\delta$関数}
  本問では、Diracの$\delta$関数が満たす基本的性質を、$\delta$関数の定義
  \begin{align*}
    \int_{-\infty}^{\infty} f(x)\delta(x-a)\,dx = f(a)
  \end{align*}
  に基づいて示す。また、Heavisideの階段関数との関係についても確認する。

  \subsubsection*{(i)(a)$\delta(-x)=\delta(x)$}
  任意の試験関数$f(x)$に対して、次の積分を考える:
  \begin{align*}
    I = \int_{-\infty}^{\infty} f(x)\delta(-x)\,dx。
  \end{align*}
  ここで変数変換
  \begin{align*}
    u=-x
  \end{align*}
  を行う。このとき
  \begin{align*}
    dx=-du
  \end{align*}
  であるから
  \begin{align*}
    I
    &= \int_{\infty}^{-\infty} f(-u)\delta(u)(-du) \\
    &= \int_{-\infty}^{\infty} f(-u)\delta(u)\,du \\
    &= f(0)。
  \end{align*}
  一方、
  \begin{align*}
    \int_{-\infty}^{\infty} f(x)\delta(x)\,dx = f(0)
  \end{align*}
  であるから、
  \begin{align*}
    \delta(-x)=\delta(x)
  \end{align*}
  が成り立つ。

  \subsubsection*{(i)(b)$\delta(ax)=\dfrac1{|a|}\delta(x)$}

  任意の試験関数$f(x)$に対して、
  \begin{align*}
    I = \int_{-\infty}^{\infty} f(x)\delta(ax)\,dx
  \end{align*}
  を考える。

  ここで変数変換
  \begin{align*}
    u=ax
  \end{align*}
  を用いる。このとき
  \begin{align*}
    dx=\frac{du}{|a|}
  \end{align*}
  であるから、
  \begin{align*}
    I
    &= \frac{1}{|a|}\int_{-\infty}^{\infty} f\left(\frac{u}{a}\right)\delta(u)\,du \\
    &= \frac{1}{|a|}f(0)。
  \end{align*}
  したがって、
  \begin{align*}
    \delta(ax)=\frac{1}{|a|}\delta(x)
  \end{align*}
  が成り立つ。

  \subsubsection*{(i)(c)$x\delta(x)=0$}

  任意の試験関数$f(x)$に対して、
  \begin{align*}
    \int_{-\infty}^{\infty} f(x)x\delta(x)\,dx
    &= \int_{-\infty}^{\infty} (xf(x))\delta(x)\,dx \\
    &= 0 \cdot f(0) \\
    &= 0。
  \end{align*}
  よって、任意の$f(x)$に対して積分値が0となるので、
  \begin{align*}
    x\delta(x)=0
  \end{align*}
  が成り立つ。

  \subsubsection*{(i)(d)$\delta(x^2-a^2)=\dfrac1{2|a|}\qty[\delta(x+a)+\delta(x-a)]$}

  一般に、関数$g(x)$が単純零点$x_i$を持つとき，
  \begin{align*}
    \delta(g(x))
    =
    \sum_i
    \frac{\delta(x-x_i)}{|g'(x_i)|}
  \end{align*}
  が成り立つ。

  ここで
  \begin{align*}
    g(x)=x^2-a^2
  \end{align*}
  とおくと，
  \begin{align*}
    g(x)=0
  \end{align*}
  の解は
  \begin{align*}
    x=\pm a
  \end{align*}
  である。また，
  \begin{align*}
    g'(x)=2x
  \end{align*}
  であるから，
  \begin{align*}
    |g'(\pm a)|=2|a|
  \end{align*}
  となる。したがって，
  \begin{align*}
    \delta(x^2-a^2)
    &= \frac{\delta(x-a)}{|2a|}
    + \frac{\delta(x+a)}{|{-2a}|} \\
    &= \frac{1}{2|a|}\qty[\delta(x+a)+\delta(x-a)]。
  \end{align*}

  \subsubsection*{(i)(e)$\int_{-\infty}^{\infty}\delta(a-x)\delta(x-b)\,dx=\delta(a-b)$}

  $\delta$関数の偶性
  \begin{align*}
    \delta(a-x)=\delta(x-a)
  \end{align*}
  を用いると、
  \begin{align*}
    \int_{-\infty}^{\infty}\delta(a-x)\delta(x-b)\,dx
    &= \int_{-\infty}^{\infty}\delta(x-a)\delta(x-b)\,dx \\
    &= \delta(a-b)。
  \end{align*}
  ここで最後の等号では，
  \begin{align*}
    f(x)=\delta(x-a)
  \end{align*}
  として$\delta$関数の定義を用いた。

  \subsubsection*{(i)(f)$x\delta'(x)=-\delta(x)$}

  $\delta$関数の微分は分布の意味で定義され、
  \begin{align*}
    \int_{-\infty}^{\infty} f(x)\delta'(x)\,dx
    = -f'(0)
  \end{align*}
  が成り立つ。

  ここで、
  \begin{align*}
    g(x)=xf(x)
  \end{align*}
  とおく。このとき，
  \begin{align*}
    g'(x)=f(x)+xf'(x)
  \end{align*}
  であるから，
  \begin{align*}
    g'(0)=f(0)
  \end{align*}
  となる。

  以上を用いると、
  \begin{align*}
    \int_{-\infty}^{\infty} f(x)x\delta'(x)\,dx
    &= \int_{-\infty}^{\infty} g(x)\delta'(x)\,dx \\
    &= -g'(0) \\
    &= -f(0) \\
    &= \int_{-\infty}^{\infty} f(x)(-\delta(x))\,dx。
  \end{align*}
  したがって、
  \begin{align*}
    x\delta'(x)=-\delta(x)
  \end{align*}
  が成り立つ。

  \subsubsection*{(ii)$\bbn{\theta(x)}{x}=\delta(x)$}

  Heavisideの階段関数
  \begin{align*}
    \theta(x)
    =
    \begin{cases}
      0 &(x\le0) \\
      1 &(x>0)
    \end{cases}
  \end{align*}
  を考える。

  任意の試験関数$f(x)$に対して、
  \begin{align*}
    \int_{-\infty}^{\infty} f(x)\theta'(x)\,dx
  \end{align*}
  を計算する。分布の微分の定義より、
  \begin{align*}
    \int_{-\infty}^{\infty} f(x)\theta'(x)\,dx
    &= -\int_{-\infty}^{\infty} f'(x)\theta(x)\,dx \\
    &= -\int_0^{\infty} f'(x)\,dx \\
    &= -[f(x)]_0^{\infty}。
  \end{align*}
  試験関数$f(x)$は無限遠方で0に収束するので、
  \begin{align*}
    f(\infty)=0
  \end{align*}
  である。したがって，
  \begin{align*}
    -[f(x)]_0^{\infty}
    &= -(0-f(0)) \\
    &= f(0)。
  \end{align*}
  一方、
  \begin{align*}
    \int_{-\infty}^{\infty} f(x)\delta(x)\,dx=f(0)
  \end{align*}
  であるから、
  \begin{align*}
    \bbn{\theta(x)}{x}=\delta(x)
  \end{align*}
  が示された。
\end{document}